\documentclass[12pt]{article}
\usepackage[a4paper,margin=1in]{geometry}
\usepackage[utf8]{inputenc}
\usepackage[T1]{fontenc}
\usepackage{amsmath,amssymb}
\usepackage{textcomp}
\usepackage{parskip}
\begin{document}
\begin{center}
{\Large \textbf{2011 STEP I}}
\end{center}

\subsection*{Page 1}
\begin{flushleft}
\bigskip\textbf{Section A: Pure Mathematics}\par
a b ay\par
\bigskip\textbf{Question 1.} (i) Show that the gradient of the curve + = 1, where b = 0, is .\par
\[x y 6 - bx2\]
a b\par
\[The point (p, q) lies on both the straight line ax + by = 1 and the curve + = 1 ,\]
x y\par
\[where ab = 0. Given that, at this point, the line and the curve have the same gradient,\]
\[show that p = q .\]
±\par
2 2\par
\[Show further that either (a b) = 1 or (a + b) = 1 .\]
-\par
\textbf{(ii)} Show that if the straight line ax + by = 1, where ab = 0, is a normal to the curve\par
a b\par
2 2 1\par
\[= 1, then a b = .\]
x - y - 2\par
1 x\par
e\par
\bigskip\textbf{Question 2.} The number E is defined by E = dx .\par
\[Z^0 1 + x\]
Show that\par
1 x\par
xe\par
\[dx = e 1 E ,\]
\[Z^0 1 + x - -\]
1 2 x\par
x e\par
and evaluate dx in terms of e and E.\par
\[Z^0 1 + x\]
Evaluate also, in terms of E and e as appropriate:\par
1 1-x\par
e 1+x\par
\textbf{(i)} dx ;\par
\[Z^0 1 + x\]
\[\sqrt\{2\} x^2\]
e\par
\textbf{(ii)} dx .\par
\[Z^1 x\]
\[9465 Jun11\]
\end{flushleft}
\subsection*{Page 2}
\begin{flushleft}
\bigskip\textbf{Question 3.} Prove the identity\par
1 1\par
\[4 sin θ sin( \pi θ) sin( \pi + θ) = sin 3θ . ( )\]
3 3\par
- ∗\par
\textbf{(i)} By differentiating ( ), or otherwise, show that\par
∗\par
\[cot 1 \pi cot 2 \pi + cot 4 \pi = \sqrt\{3\} .\]
9 9 9\par
-\par
\textbf{(ii)} By setting θ = \pi φ in ( ), or otherwise, obtain a similar identity for cos 3θ and deduce\par
- ∗\par
that\par
1 1\par
\[cot θ cot( \pi θ) cot( \pi + θ) = cot 3θ .\]
3 3\par
-\par
Show that\par
\[cosec 1 \pi cosec 5 \pi + cosec 7 \pi = 2\sqrt\{3\} .\]
9 9 9\par
-\par
2 2\par
\bigskip\textbf{Question 4.} The distinct points P and Q, with coordinates (ap , 2ap) and (aq , 2aq) respectively, lie on\par
\[the curve y = 4ax. The tangents to the curve at P and Q meet at the point T . Show that\]
\[T has coordinates apq, a(p + q) . You may assume that p = 0 and q = 0.\]
6 6\par
The point F has coordinates (a, 0) and φ is the angle T F P . Show that\par
pq + 1\par
\[cos φ =\]
\[(p^2 + 1)(q^2 + 1)\]
p\par
and deduce that the line F T bisects the angle P F Q.\par
\bigskip\textbf{Question 5.} Given that 0 < k < 1, show with the help of a sketch that the equation\par
\[sin x = kx ( )\]
∗\par
\[has a unique solution in the range 0 < x < \pi.\]
Let\par
\[\pi\]
\[I = sin x kx dx .\]
\[Z^0 -\]
\{ \{\par
\{ \{\par
Show that\par
\[\pi sin α\]
\[I = 2 cos α α sin α ,\]
2α - -\par
where α is the unique solution of ( ).\par
∗\par
Show that I, regarded as a function of α, has a unique stationary value and that this stationary\par
value is a minimum. Deduce that the smallest value of I is\par
\[\pi\]
2 cos .\par
\[- \sqrt\{2\}\]
\[9465 Jun11 [Turn over\]
\end{flushleft}
\subsection*{Page 3}
\begin{flushleft}
r 3\par
\bigskip\textbf{Question 6.} Use the binomial expansion to show that the coefficient of x in the expansion of (1 x)\par
-\par
1 -\par
is (r + 1)(r + 2) .\par
r\par
\textbf{(i)} Show that the coefficient of x in the expansion of\par
1 x + 2x\par
-\par
(1 x)3\par
-\par
is r + 1 and hence find the sum of the series\par
2 5 10 17 26 37 50\par
1 + + + + + + + + .\par
2 4 8 16 32 64 128 \cdots\par
\textbf{(ii)} Find the sum of the series\par
9 25 9 49\par
1 + 2 + + 2 + + + + .\par
4 16 8 64 \cdots\par
1 2\par
\bigskip\textbf{Question 7.} In this question, you may assume that ln(1 + x) x x when x is small.\par
\approx - | |\par
The height of the water in a tank at time t is h. The initial height of the water is H and water\par
flows into the tank at a constant rate. The cross-sectional area of the tank is constant.\par
\textbf{(i)} Suppose that water leaks out at a rate proportional to the height of the water in the\par
tank, and that when the height reaches α H, where α is a constant greater than 1, the\par
height remains constant. Show that\par
dh\par
\[= k(α H h) ,\]
dt -\par
for some positive constant k. Deduce that the time T taken for the water to reach height\par
αH is given by\par
\[kT = ln 1 + ,\]
\lambda α \mu\par
and that kT α for large values of α.\par
-\par
\approx\par
\textbf{(ii)} Suppose that the rate at which water leaks out of the tank is proportional to \sqrth (instead\par
of h), and that when the height reaches α H, where α is a constant greater than 1, the\par
height remains constant. Show that the time T taken for the water to reach height αH\par
is given by\par
\[cT = 2\sqrt\{H\} 1 \sqrtα + α ln 1 +\]
\[\lambda - \lambda \sqrtα \mu\mu\]
\[for some positive constant c, and that cT \sqrt\{H\} for large values of α.\]
\approx\par
\[9465 Jun11\]
\end{flushleft}
\subsection*{Page 4}
\begin{flushleft}
\bigskip\textbf{Question 8.} (i) The numbers m and n satisfy\par
3 3 2\par
\[m = n + n + 1 . ( )\]
∗\par
\[(a) Show that m > n. Show also that m < n + 1 if and only if 2n + 3n > 0 . Deduce\]
\[that n < m < n + 1 unless 3 6 n 6 0 .\]
-\par
(b) Hence show that the only solutions of ( ) for which both m and n are integers are\par
∗\par
\[(m, n) = (1, 0) and (m, n) = (1, 1).\]
-\par
\textbf{(ii)} Find all integer solutions of the equation\par
3 3 2\par
\[p = q + 2q 1 .\]
-\par
\[9465 Jun11 [Turn over\]
\end{flushleft}
\subsection*{Page 5}
\begin{flushleft}
\bigskip\textbf{Section B: Mechanics}\par
\bigskip\textbf{Question 9.} A particle is projected at an angle θ above the horizontal from a point on a horizontal plane.\par
The particle just passes over two walls that are at horizontal distances d1 and d2 from the\par
\[point of projection and are of heights d^2 and d^1, respectively. Show that\]
2 2\par
d + d d + d\par
1 1 2 2\par
\[tan θ = .\]
d d\par
1 2\par
Find (and simplify) an expression in terms of d1 and d2 only for the range of the particle.\par
\bigskip\textbf{Question 10.} A particle, A, is dropped from a point P which is at a height h above a horizontal plane.\par
A second particle, B, is dropped from P and first collides with A after A has bounced on the\par
plane and before A reaches P again. The bounce and the collision are both perfectly elastic.\par
Explain why the speeds of A and B immediately before the first collision are the same.\par
\[The masses of A and B are M and m, respectively, where M > 3m, and the speed of the\]
particles immediately before the first collision is u. Show that both particles move upwards\par
after their first collision and that the maximum height of B above the plane after the first\par
collision and before the second collision is\par
4M (M m)u\par
h + - .\par
(M + m)2g\par
\bigskip\textbf{Question 11.} A thin non-uniform bar AB of length 7d has centre of mass at a point G, where AG = 3d.\par
A light inextensible string has one end attached to A and the other end attached to B. The\par
string is hung over a smooth peg P and the bar hangs freely in equilibrium with B lower\par
than A. Show that\par
\[3 sin α = 4 sin β ,\]
where α and β are the angles P AB and P BA, respectively.\par
Given that cos β = and that α is acute, find in terms of d the length of the string and show\par
that the angle of inclination of the bar to the horizontal is arctan .\par
\[9465 Jun11\]
\end{flushleft}
\subsection*{Page 6}
\begin{flushleft}
\bigskip\textbf{Section C: Probability and Statistics}\par
\bigskip\textbf{Question 12.} I am selling raffle tickets for \pounds 1 per ticket. In the queue for tickets, there are m people each\par
with a single \pounds 1 coin and n people each with a single \pounds 2 coin. Each person in the queue wants\par
to buy a single raffle ticket and each arrangement of people in the queue is equally likely to\par
occur. Initially, I have no coins and a large supply of tickets. I stop selling tickets if I cannot\par
give the required change.\par
\textbf{(i)} In the case n = 1 and m > 1, find the probability that I am able to sell one ticket to\par
each person in the queue.\par
\textbf{(ii)} By considering the first three people in the queue, show that the probability that I am\par
m 1\par
\[able to sell one ticket to each person in the queue in the case n = 2 and m > 2 is - .\]
m + 1\par
\textbf{(iii)} Show that the probability that I am able to sell one ticket to each person in the queue\par
m 2\par
\[in the case n = 3 and m > 3 is - .\]
m + 1\par
\bigskip\textbf{Question 13.} In this question, you may use without proof the following result:\par
\[4 x^2 dx = 2 arcsin( 1 x) + 1 x 4 x^2 + c .\]
2 2\par
Z - -\par
p p\par
A random variable X has probability density function f given by\par
\[2k a 6 x < 0\]
-\par
\par
\[f(x) =  k\sqrt\{4\} x^2 0 6 x 6 2\]
\par
  -\par
0 otherwise,\par
\par
\par
\par
\par
\par
where k and a are positive constants.\par
\textbf{(i)} Find, in terms of a, the mean of X.\par
\textbf{(ii)} Let d be the value of X such that P(X > d) = . Show that d < 0 if 2a > 9\pi and find\par
an expression for d in terms of a in this case.\par
\textbf{(iii)} Given that d = \sqrt2, find a.\par
\[9465 Jun11 [Turn over\]
\end{flushleft}
\end{document}
